\section{Method}

\subsection{Overview}

The proposed method is denoted as EA-RG-MAPPO-S, namely Edge-Aware Role Graph MAPPO with Staged random-radius fine-tuning. It contains a local-observation encoder, an edge-aware role graph encoder, a centralized critic, and a staged random-radius fine-tuning procedure. The actor input is the concatenation of the local observation representation and the graph representation, while the critic uses the centralized global state. The method overview is shown in Fig.~\ref{fig:method_overview_en}.

\begin{figure*}[t]
    \centering
    \includegraphics[width=0.92\textwidth]{method_overview_ea_rg_mappo_s.png}
    \caption{Overview of EA-RG-MAPPO-S.}
    \label{fig:method_overview_en}
\end{figure*}

\subsection{Role Graph Construction}

At each time step, a dynamic graph $G^t=(V^t,E^t)$ is constructed. The node set includes pursuer nodes and the target node. Each node contains normalized position, heading, speed, and platform-related state features. A role indicator is used to distinguish UAV nodes from target nodes. The node feature and role embedding are concatenated before entering the graph encoder:
\begin{equation}
h_i^0 = f_{\mathrm{in}}\left([x_i,\mathrm{Emb}(\mathrm{role}_i)]\right).
\end{equation}
The role graph allows the model to distinguish teammates from the target, avoiding a homogeneous treatment of all graph nodes.

\subsection{Relative Edge-Feature Enhanced Attention}

For the directed edge from node $i$ to node $j$, a relative edge feature vector $e_{ij}$ is constructed:
\begin{equation}
e_{ij}=[\Delta x_{ij},\Delta y_{ij},d_{ij}/d_{\max},d_{ij}/R_c,
\cos\beta_{ij},\sin\beta_{ij},\Delta v^x_{ij},\Delta v^y_{ij},
\mathbb{I}_{ij}^{\mathrm{comm}},\mathbb{I}_{j}^{\mathrm{target}}].
\end{equation}
Here $d_{ij}$ is the distance between nodes, $d_{\max}$ is the environment-scale normalization constant, $\beta_{ij}$ is the relative bearing, $\mathbb{I}_{ij}^{\mathrm{comm}}$ indicates communication reachability, and $\mathbb{I}_{j}^{\mathrm{target}}$ indicates whether node $j$ is the target. The edge-aware attention score is computed as:
\begin{equation}
\mathrm{score}_{ij}=\mathrm{LeakyReLU}\left(a^\top[Wh_i,Wh_j]\right)+g_e(e_{ij}),
\end{equation}
where $g_e$ is an edge-feature scoring network. A masked softmax is then applied over reachable neighbors:
\begin{equation}
\alpha_{ij}=\frac{\exp(\mathrm{score}_{ij})}{\sum_{k\in\mathcal{N}_i}\exp(\mathrm{score}_{ik})}, \quad j\in\mathcal{N}_i.
\end{equation}
The node representation is updated by:
\begin{equation}
h_i'=\tanh\left(\sum_{j\in\mathcal{N}_i}\alpha_{ij}Wh_j\right).
\end{equation}
Compared with standard GAT, which computes attention only from node embeddings, EA-RG-MAPPO-S directly uses distance, bearing, velocity difference, and communication reachability to adjust neighbor weights.

\subsection{MAPPO Optimization}

The policy network outputs the action distribution $\pi_\theta(a_i|o_i,G)$ for each agent. MAPPO is used with the clipped surrogate objective. The probability ratio is:
\begin{equation}
\rho_i^t(\theta)=\frac{\pi_\theta(a_i^t|o_i^t,G^t)}{\pi_{\theta_{\mathrm{old}}}(a_i^t|o_i^t,G^t)}.
\end{equation}
The policy loss is:
\begin{equation}
L_{\mathrm{policy}}=-\mathbb{E}\left[\min\left(\rho_i^t A_i^t,\mathrm{clip}(\rho_i^t,1-\epsilon,1+\epsilon)A_i^t\right)\right].
\end{equation}
The total loss is:
\begin{equation}
L=L_{\mathrm{policy}}+c_vL_{\mathrm{value}}-c_HL_{\mathrm{entropy}}+c_{\mathrm{aux}}L_{\mathrm{aux}}.
\end{equation}
$L_{\mathrm{aux}}$ is an optional auxiliary loss. Since the current target-intent branch does not achieve reliable balanced accuracy, the main conclusions do not depend on that auxiliary branch.

\subsection{Staged Random-Radius Fine-Tuning}

Training only under a fixed communication radius may make the policy over-adapt to a single communication topology. This paper uses a two-stage training strategy. In the first stage, the edge-aware role graph policy is trained under a fixed communication radius $R_c=8$ to learn basic coordination. In the second stage, the policy is initialized from the first-stage checkpoint and fine-tuned with a smaller learning rate. For each episode, the communication radius is sampled from:
\begin{equation}
R_c\sim U(4,10).
\end{equation}
The purpose of this stage is to improve adaptation to different communication topologies rather than merely increasing the total training time.
